\par\Needspace{7\baselineskip}\bigskip
\begin{center}
\pdfbookmark[2]{Encke an Gauss, Berlin, 4. October 1834}{pallas-brief-18341004}
\textsc{Encke} an \textsc{Gauss}, Berlin, 4. October 1834.
\end{center}
\noindent\hspace*{\parindent}\{. . . . .\quad Die Methode der partiellen Störungen, welche Sie mir im Jahr 1811 vorzutragen die Güte hatten, hatten Sie mir damals unter der Äusserung mitgetheilt, dass Sie sie nicht verbreitet zu sehen wünschten, da Sie selbst etwas darüber mittheilen wollten. Ich habe sie seitdem beständig angewandt und bin jetzt durch das Jahrbuch genöthigt, sie auf die vier kleinen Planeten und ausserdem noch auf den Cometen von kurzer Umlaufszeit fortwährend anzuwenden. \textsc{Heiligensteins} Tod hat es auch für die Ceres nothwendig gemacht, was in diesen Ferien geschehen ist. Indessen sehe ich voraus, dass es mir in Zukunft allein nicht mehr möglich sein wird, und dass, selbst wenn ein oder der andere Ihrer Schüler einen oder den andern Planeten übernehmen wollte, mir damit nicht geholfen sein würde, weil ich höchst wahrscheinlich selbst die Arbeiten über den Cometen einem Andern zu übertragen genöthigt sein werde. Ich möchte Sie deswegen ersuchen, mir zu erlauben, diese Mittheilung Ihrer Methoden Jedem machen zu dürfen, dem ich eine solche Arbeit anvertraue, da ohne die freie Mittheilung eine Hülfe nicht möglich ist und noch mehr würden Sie mich beglücken, wenn Sie mir erlauben wollten, etwa in einem Anhange des Jahrbuchs die Methode ausführlich vorzutragen.\quad\mbox{. . . . .\}}
\par
